\subsubsection*{1.2.7.}
Let $D$ be a disc, let $0\in D$, put $D^*=D-\{0\}$, let
$f:X\longrightarrow D$ be a continuous map, and put
$X_0=f^{-1}(0)$.  We shall define a functor
\[
\Psi:(\text{sheaves on }X)
\longrightarrow(\text{sheaves on }X_0\times\mathcal D).
\]
Let $\widetilde D^*$ be a universal covering of $D^*$, let
$\widetilde D$ be as in 1.1.9, and put
\[
\overline X=X\times_D\widetilde D,
\qquad
X^*=X-X_0=X\times_DD^*,
\qquad
\overline X^*=\overline X-X_0=X\times_D\widetilde D^*.
\]
There is a commutative diagram
\begin{equation}\tag{1.2.7.1}
\begin{tikzcd}
X_0 \arrow[r, hook, "\overline i"] \arrow[d, equal]
  & \overline X \arrow[d, "p"]
  & \overline X^* \arrow[l, hook, "\overline j"']
      \arrow[d, "p'"] \\
X_0 \arrow[r, hook, "i"]
  & X
  & X^* \arrow[l, hook, "j"']
\end{tikzcd}
\end{equation}
cartesian over the diagram
\begin{equation}\tag{1.2.7.2}
\begin{tikzcd}
\{0\} \arrow[r, hook, "\overline i"] \arrow[d, equal]
  & \widetilde D \arrow[d, "p"]
  & \widetilde D^* \arrow[l, hook, "\overline j"']
      \arrow[d, "p'"] \\
\{0\} \arrow[r, hook, "i"]
  & D
  & D^* \arrow[l, hook, "j"']
\end{tikzcd}
\end{equation}
\noindent\emph{Editorial note.}  In (1.2.7.1) the printed top-right
node lacks the overline.  The preceding definition and the map $p'$
force that node to be $\overline X^*$, which is the notation used
here.

For a sheaf $F$ on $X$, we set (with the notation of 1.2.4(c))
\[
\Psi(F)=
\bigl(i^*F,\ \overline i^*\overline j_*((jp')^*F),\ \alpha\bigr),
\]
where
\[
\begin{split}
\alpha:i^*F=\overline i^*p^*F
&\longrightarrow
\overline i^*\overline j_*((jp')^*F)\\
&=\overline i^*\overline j_*\overline j^*p^*F
\end{split}
\]
is obtained, after applying $\overline i^*$, from the adjunction map
$\mathrm{Id}\longrightarrow\overline j_*\overline j^*$.

We also define a functor
\[
\Psi_\eta:(\text{sheaves on }X^*)
\longrightarrow(\text{sheaves on }X_0\times\mathcal D^*)
\]
by
\[
\Psi_\eta(F)=\overline i^*\overline j_*({p'}^*F).
\]

\subsubsection*{1.2.8.}
The functoriality properties of $\Psi$ are the same as those of its
analogue studied in XIII, 1.3.5--1.3.9.

